\subsection{FUNCTORIALITY OF EXPONENTIAL MAPPINGS}

PROPOSITION 8. Let G and H be Lie group germs, h a morphism of G into H and \(\phi_{\mathrm{G}}\) and \(\phi_{\mathrm{H}}\) exponential mappings of G and H. There exists a neighbourhood V of 0 in L(G) such that \(h \circ \phi_{\mathrm{G}}\) and \(\phi_{\mathrm{H}} \circ \mathrm{L}(h)\) coincide on V.

Let \(\mathrm{L}(\mathrm{G})\) and \(\mathrm{L}(\mathrm{H})\) be given norms defining their topologies and such that \(\| [x,y]\| \leqslant \| x\| \| y\|\) for all \(x\) and \(y\) . It can be assumed that \(\mathrm{G}\) (resp. H) is the Lie group germ defined by \(\mathrm{L}(\mathrm{G})\) (resp. \(\mathrm{L}(\mathrm{H})\) ), so that \(\phi_{\mathrm{G}}\) (resp. \(\phi_{\mathrm{H}}\) ) coincides with \(\mathrm{Id}_{\mathrm{G}}\) (resp. \(\mathrm{Id}_{\mathrm{H}}\) ) in a neighbourhood of 0. On the other hand, there exists a symmetric open neighbourhood W of 0 in \(\mathrm{L}(\mathrm{G})\) such that \(\mathrm{L}(h)\) is a morphism of the Lie group germ W into H. By Theorem 1, \(\mathrm{L}(h)\) coincides with \(h\) on a neighbourhood of 0, whence the proposition.

Loosely speaking, if G and H are identified in a neighbourhood of the identity element with L(G) and L(H) by means of exponential mappings, every morphism of G into H is linear in a neighbourhood of 0.

COROLLARY 1. Let G be a Lie group germ, G' a Lie subgroup germ of G and \(\phi\) an exponential mapping of G.

\begin{enumerate}
\def\labelenumi{(\roman{enumi})}
\item
  There exists an open neighbourhood \(\mathbf{V}\) of 0 in \(\mathrm{L}(\mathbf{G}')\) such that \(\phi | \mathbf{V}\) is an isomorphism of the manifold \(\mathbf{V}\) onto an open neighbourhood of \(e\) in \(\mathbf{G}'\) .
\item
  Let \(x \in \mathrm{L}(\mathrm{G})\) . The following conditions are equivalent: (a) \(x \in \mathrm{L}(\mathrm{G}')\) ; (b) \(\phi(\lambda x) \in \mathrm{G}'\) for \(|\lambda|\) sufficiently small.
\item
  is obtained by applying Proposition 8 to the canonical injection of \(\mathbf{G}'\) into \(\mathbf{G}\) and (ii) follows from (i).
\end{enumerate}

COROLLARY 2. Let G be a Lie group, \(\rho\) an analytic linear representation of G and \(\phi\) an exponential mapping of G. There exists a neighbourhood V of 0 in L(G) such that

for all \(x \in V\) .

\[
\rho (\phi (x)) = \exp (\mathrm{L} (\rho) x)
\]

This follows from Proposition 8 and Example 2 of no. 3.

COROLLARY 3. Let G be a Lie group and \(\phi\) an exponential mapping of G.

\begin{enumerate}
\def\labelenumi{(\roman{enumi})}
\tightlist
\item
  There exists a neighbourhood V of 0 in L(G) such that
\end{enumerate}

\[
\operatorname{Ad} (\phi (x)) = \exp \operatorname{ad} x.
\]

for all \(x \in V\) .

\begin{enumerate}
\def\labelenumi{(\roman{enumi})}
\setcounter{enumi}{1}
\tightlist
\item
  If \(g \in \mathbf{G}\) , there exists a neighbourhood \(\mathbf{W}\) of 0 in \(\mathrm{L}(\mathbf{G})\) such that
\end{enumerate}

\[
g \phi (x) g ^ {- 1} = (\mathrm{Ad} g. x)
\]

for all \(x \in W\) .

\begin{enumerate}
\def\labelenumi{(\roman{enumi})}
\tightlist
\item
  follows from Corollary 2 and § 3, no. 12, Proposition 44.
\end{enumerate}

\[
g.
\]

